\chapter{Introducción}

\section{Sobre las notas}
Estas notas explican recordatorios y cosas basicas en programación competitiva. 
Fueron escritas con la idea de usarla como notebook en competencias. 



\section{Cosas de C++}

\subsection{Tamaño de los tipos de datos}
\begin{itemize}
	\item \textbf{Tamaño de un long long}: Es un número de 64 bits, \( 2^{64} \) o aproximadamente \( 10^{19} \).
\end{itemize}

\subsection{Lectura de líneas con espacios en C++}

En C++, para leer una línea completa que puede contener espacios (como una oración o una entrada con varios números separados), se utiliza:

\begin{lstlisting}
	getline(std::cin, s);
\end{lstlisting}

Sin embargo, si antes se usó 

\begin{lstlisting}[language=C++]
	std::cin >> variable;
\end{lstlisting}

entonces \texttt{getline()} puede leer una línea vacía. Para evitarlo, se debe limpiar el buffer con:

\begin{lstlisting}[language=C++]
	std::cin.ignore();
\end{lstlisting}



\begin{lstlisting}[language=C++]
	cin.ignore();
\end{lstlisting}



\subsection{Conversiones }

\subsubsection{String a Numero y de Numero a String }
\begin{itemize}
	\item \textbf{De string a int}: \texttt{int x = stoi("123");}
	\item \textbf{De string a long long}: \texttt{ll y = stoll("9876543210");} 
	\item \textbf{De numero a string}: \texttt{string s = to\_string(x);}
\end{itemize}

\subsubsection{Conversiones de char a número posicionales }

\paragraph{De char numeral a número:}
Si el \texttt{char} representa un dígito (por ejemplo, \texttt{'0'} a \texttt{'9'}), se puede obtener el valor numérico restando el carácter \texttt{'0'}:

\begin{lstlisting}[language=C++]
	ll x = s[0] - '0'; // convierte '0' a 0, '5' a 5, etc.
\end{lstlisting}

\paragraph{De char minúscula a número ($a \to 1$, $b \to 2$, $c \to 3$, ...):}
Dado un carácter en minúscula \( c \in \{a, b, \dots, z\} \), se puede obtener su posición en el alfabeto con:

\[
\text{posición}(c) = c - 'a' + 1
\]

Ejemplos:
\begin{lstlisting}[language=C++]
	int pos_b = 'b' - 'a' + 1; // pos_b == 2
	int pos_z = 'z' - 'a' + 1; // pos_z == 26
\end{lstlisting}

\paragraph{De char mayúscula a número ($A \to 1$, $B \to 2$, $C \to 3$, ...):}
Para caracteres en mayúscula \( C \in \{A, B, \dots, Z\} \), se usa la misma lógica:

\[
\text{posición}(C) = C - 'A' + 1
\]

Ejemplos:
\begin{lstlisting}[language=C++]
	int pos_B = 'B' - 'A' + 1; // pos_B == 2
	int pos_Z = 'Z' - 'A' + 1; // pos_Z == 26
\end{lstlisting}

\subsubsection{De mayuscula a minuscula y visceversa}


\begin{itemize}
	\item \textbf{Convertir un carácter a mayúscula:}
	\begin{lstlisting}[language=C++]
		char c = 'b';
		c = toupper(c); // c ahora es 'B'
	\end{lstlisting}
	
	\item \textbf{Convertir un carácter a minúscula:}
	\begin{lstlisting}[language=C++]
		char c = 'G';
		c = tolower(c); // c ahora es 'g'
	\end{lstlisting}
\end{itemize}

\subsubsection{Conversiones a binarias y visceversas}

\paragraph{Convertir número a string binario:}
Para convertir un número entero a su representación binaria en forma de string, se puede usar la función:

\begin{lstlisting}[language=C++]
	
	ll num = 13;
	string binary = bitset<64>(num).to_string(); // "0000...1101"
	
	// O sin ceros a la izquierda:
	binary = bitset<64>(num).to_string();
	binary.erase(0, binary.find('1')); // remueve ceros iniciales
\end{lstlisting}

\paragraph{Convertir string binario a número:}
Para convertir un string que representa un número binario a su valor decimal, se usa \texttt{stoll()} con base 2:

\begin{lstlisting}[language=C++]
	string binary = "1101";
	ll num = stoll(binary, nullptr, 2); // num == 13
	
	// Otros ejemplos:
	ll a = stoll("0101", nullptr, 2);  // a == 5
	ll b = stoll("11111111", nullptr, 2); // b == 255
\end{lstlisting}

\textbf{Nota:} El parámetro \texttt{nullptr} indica que no se necesita guardar la posición del primer carácter no convertido, y el \texttt{2} especifica que se usa base binaria.

\subsection*{Ordenar partes (subrangos) de un vector}

En C++, se puede ordenar una parte de un vector (es decir, un subrango) utilizando la función \texttt{sort} de la STL, especificando iteradores de inicio y fin.

Esto es especialmente útil cuando se necesita dividir un conjunto en mitades o grupos, por ejemplo, para aplicar diferentes criterios de ordenamiento en cada parte.

\begin{lstlisting}[language=C++]
	// Supongamos que tenemos un vector de arreglos de 3 enteros:
	vector<array<ll, 3>> v(n);
	
	// Ordenamos solo la primera mitad del vector por el segundo campo (indice 1):
	sort(v.begin(), v.begin() + n/2, [](const array<ll, 3> &a, const array<ll, 3> &b) {
		return a[1] < b[1];
	});
	
	// Ordenamos la segunda mitad tambien por el segundo campo:
	sort(v.begin() + n/2, v.end(), [](const array<ll, 3> &a, const array<ll, 3> &b) {
		return a[1] < b[1];
	});
\end{lstlisting}

\textbf{Aplicación:} esta técnica aparece frecuentemente en problemas donde hay que emparejar elementos de dos mitades, o donde diferentes partes del conjunto deben ser tratadas con diferentes estrategias.

\section{Cosas comunes en programación competitiva}


\subsection{Notación BIG O}
\begin{table}[h!]
	\centering
	\begin{tabular}{|c|c|}
		\hline
		\textbf{Tamaño del input: $N$} & \textbf{Peor complejidad aceptada} \\ \hline
		$N \leq 22 \ (\approx 2 \cdot 10^1)$ & $O(2^N)$ \\ \hline
		$N \leq 100 \ (10^2)$ & $O(N^3)$ \\ \hline
		$N \leq 1000 \ (10^3)$ & $O(N^2)$ \\ \hline
		$N \leq 1{,}000{,}000 \ (10^6)$ & $O(N \log N)$ \\ \hline
		$N \leq 1{,}000{,}000{,}000 \ (10^9)$ & $O(N),\ O(\log N),\ O(1)$ \\ \hline
	\end{tabular}
	\caption{Relación entre tamaño de input y la peor complejidad aceptada.}
\end{table}

\subsection{Resize() en c++}
La funci\'on \texttt{std::vector::resize()} ajusta el n\'umero de elementos en un vector. Su acci\'on depende del tama\~no deseado:

\begin{itemize}
	\item \textbf{Si el nuevo tama\~no es mayor que el actual:} Se a\~naden nuevos elementos al final. Opcionalmente, se puede especificar un valor de inicializaci\'on: \texttt{mi\_vector.resize(10, -1)}.
	\item \textbf{Si el nuevo tama\~no es menor que el actual:} Se eliminan los elementos del final del vector hasta alcanzar el tama\~no especificado.
\end{itemize}
\begin{lstlisting}[style=mycpp, caption=Uso eficiente de resize]

	int main() {
		int n; cin>>n;
		vector<int> numeros;
		numeros.resize(n); // Pre-asigna memoria para n elementos
		for (int i = 0; i < n; ++i) {
			cin >> numeros[i]; // Acceso directo por \'indice
		}
		return 0;
	}
\end{lstlisting}

